\uuid{307T}
\titre{Carré d'une variable normale centrée réduite}
\auteur{Maxime Nguyen}
\organisation{AMSCC}
\datecreate{2026-11-27}
\niveau{L2}
\module{Probabilité}
\chapitre{Probabilité continue}
\sousChapitre{Loi normale}
\difficulte{2}
\theme{loi normale centrée réduite, moment d'ordre 2, changement de variable}

\contenu{
\begin{question}
Soit $Z$ une variable de loi normale centrée réduite $\mathcal{N}(0,1)$ et $\Phi$ sa fonction de répartition. Quelles affirmations sont vraies ?
\begin{answers}
\good{$\mathbb{E}(Z^2) = 1$.}
\bad{$Z^2$ suit la loi $\mathcal{N}(0,1)$.}
\bad{$\mathbb{P}(Z^2 \leq 1) = \Phi(1)$.}
\good{$\mathbb{P}(Z^2 \leq 1) = 2\Phi(1) - 1$.}
\end{answers}
\begin{explanations}
A : $\mathbb{E}(Z^2) = \mathbb{V}(Z) + \mathbb{E}(Z)^2 = 1 + 0$.

B est faux : $Z^2 \geq 0$, alors qu'une variable normale prend des valeurs négatives avec une probabilité non nulle. (Sa loi est la loi du $\chi^2$ à un degré de liberté.)

C, D : $\{Z^2 \leq 1\} = \{-1 \leq Z \leq 1\}$, de probabilité $\Phi(1) - \Phi(-1) = 2\Phi(1) - 1$.

Contre-sens visés : croire qu'une fonction d'une variable normale est
normale ; résoudre $Z^2 \leq 1$ en $Z \leq 1$ (C).
\end{explanations}
\end{question}
}
